\documentclass[10pt,a4paper]{amsart}
\usepackage[margin=19mm]{geometry}
\usepackage{amsmath,amssymb,mathtools}
\usepackage[scr=rsfs,cal=euler]{mathalfa}
\usepackage{xfrac}
\usepackage[unicode,hidelinks,bookmarks=false]{hyperref}
\newcommand{\ie}{i.e.\ }
\newcommand{\cf}{cf.\ }
\newcommand{\resp}{respectively\ }
\newcommand{\Subsection}{\subsection}
\providecommand{\nobreakdash}{\nobreak\hbox{-}}
\setlength{\emergencystretch}{2em}
\multlinegap0pt
\usepackage{amssymb}
\newtheorem{lemma}{Lemma}
\title{Galois theory for finite fields}
\author{Askold Khovanskii}
\date{Source: arXiv:2604.08640v1 (CC BY 4.0)}
\begin{document}
\maketitle

\noindent\emph{Source and licence.} Galois theory for finite fields. Version arXiv:2604.08640v1 (CC BY 4.0), distributed by arXiv under the Creative Commons Attribution 4.0 International licence, http://creativecommons.org/licenses/by/4.0/, as recorded in the arXiv licence metadata for this exact version and shown on the abstract page https://arxiv.org/abs/2604.08640v1. Full licence notice: \texttt{licenses/arxiv-2604.08640-CC-BY-4.0.txt}.\par\smallskip

\noindent\textit{Editorial scope.} Counted unit: the article's lemma mod2 (source0.tex lines 185--191). The article's complete original proof of it is delivered with it (source0.tex lines 193--202, one \texttt{proof} environment of 10 lines). No mathematics is written, repaired or re-derived: the statement and the proof are the article's own lines, in the article's own notation, and no retained line is substituted or re-flowed.  No further source-local context is needed: every cross-reference of every retained line resolves inside the two delivered spans.  Nothing was omitted from any retained span, so the package carries no omission record.  No retained line contains a citation, so the wrapper carries no bibliography and no bibliography entry.  No retained line contains a figure environment, an image-inclusion command or drawing code, so no image is delivered and the package declares no asset dependency.  Editorial formatting only, in addition: an AMS \texttt{amsart} preamble that restates the article's own declarations and macros used by the retained lines, namely the environments \texttt{lemma} on separate counters, together with the standard packages those lines need.  The retained environments print in this document as Lemma 1 in order of appearance, whereas the article prints the same items with the numbering of its own full document.  The complete native source of the article as distributed in the arXiv e-print for this exact version is bundled unchanged as \texttt{sources/198164/source0.tex}.
\begin{lemma}\label{mod2}
\begin{enumerate}
 \item If $a\equiv \pm 1\mod 8$, then for any $k\geq 1$ we have $a^{2^k}\equiv 1\mod2^{k+3} $.

\item  If $a\equiv \pm 3\mod 8$, then for any $k \geq  1$ we have $a^{2^k}\equiv 2^{k+2} \mod2^{k+3} $.
\end{enumerate}
\end{lemma}

\begin{proof} One can prove this lemma by induction in $k$.

\begin{enumerate}
 \item For $k=1$ we have $a^{2^1}=a^2= (\pm 1+ 8 n)^2= 1\pm 16 n+16n^2= 1\mod 16= 1\mod 2^{1+3}$. If for some $k$ we have $a^{2^k}=1 +2^{k+3} n$, then $a^{2^{k+1}}=1+ 2\cdot 2^{k+3}n + 2^{2k+6}n^2 \equiv 1\mod 2^{(k+1)+3}$. The needed statement is proven.


\item For $k=1$ we have $a^{2^1}=a^2= (\pm 3+ 8 n)^2= 9 \pm 16\cdot 3n +16n^2=1+ 8\mod 16= 1+ 2^{1+2} \mod 2^{1+3}$. If for some $k$ we have $a^{2^k}=1   +2^{k+2} +2^{k+3} n $, then $a^{2^{k+1}}=(1+2^{k+2})^2+ 2\cdot 2^{k+3}n (1+2^{k+2})  +2^{2(k+3)}n^2 \equiv 1 +2^{k+3}\mod 2^{k+4}$.  The needed statement is proven.

\end{enumerate}
\end{proof}

\end{document}
